\section{Dataset}
\label{app:datasets_metrics}
\subsection{Point-wise Dataset}
\begin{table*}[ht]
\centering
\caption{Detailed track distribution of our point-wise dataset.}
\small
\begin{tabularx}{\textwidth}{>{\raggedright\arraybackslash}X c|>{\raggedright\arraybackslash}X c}
\toprule
\multicolumn{2}{c|}{\textbf{ICLR 2025}} & \multicolumn{2}{c}{\textbf{NeurIPS 2025}} \\
\midrule
1. foundation or frontier models, including LLMs & 15 & 1. Deep learning (e.g., architectures, generative models, optimization for deep networks, foundation models, LLMs) & 24 \\
2. generative models & 13 & 2. Theory (e.g., control theory, learning theory, algorithmic game theory) & 11 \\
3. applications to computer vision, audio, language, and other modalities & 11 & 3. Machine learning for sciences (e.g. climate, health, life sciences, physics, social sciences) & 7 \\
4. reinforcement learning & 11 & 4. Applications (e.g., vision, language, speech and audio, Creative AI) & 7 \\
5. datasets and benchmarks & 10 & 5. Probabilistic methods (e.g., variational inference, causal inference, Gaussian processes) & 6 \\
6. interpretability and explainable AI & 10 & 6. Social and economic aspects of machine learning (e.g., fairness, interpretability, human-AI interaction, privacy, safety, strategic behavior) & 6 \\
7. unsupervised, self-supervised, semi-supervised, and supervised representation learning & 10 & 7. Reinforcement learning (e.g., decision and control, planning, hierarchical RL, robotics) & 5 \\
8. applications to physical sciences (physics, chemistry, biology, etc.) & 7 & 8. Optimization (e.g., convex and non-convex, stochastic, robust) & 4 \\
9. learning theory & 7 & 9. Neuroscience and cognitive science (e.g., neural coding, brain-computer interfaces) & 3 \\
10. applications to robotics, autonomy, planning & 6 & 10. General machine learning (supervised, unsupervised, online, active, etc.) & 3 \\
11. learning on graphs and other geometries \& topologies & 6 & 11. Infrastructure (e.g., libraries, improved implementation and scalability, distributed solutions) & 2 \\
12. alignment, fairness, safety, privacy, and societal considerations & 6 & 12. Others & 2 \\
13. learning on time series and dynamical systems & 5 & 13. Evaluation (e.g., methodology, meta studies, replicability and validity, human-in-the-loop) & 1 \\
14. other topics in machine learning (i.e., none of the above) & 5 \\
15. optimization & 4 \\
16. applications to neuroscience \& cognitive science & 4 \\
17. neurosymbolic \& hybrid AI systems (physics-informed, logic \& formal reasoning, etc.) & 3 \\
18. causal reasoning & 1 \\
19. infrastructure, software libraries, hardware, systems, etc. & 1 \\
20. transfer learning, meta learning, and lifelong learning & 1 \\
\bottomrule
\end{tabularx}
\label{tab:track_distribution}
\end{table*}
\paragraph{Dataset Construction.}
We collect papers via the official OpenReview API\footnote{\url{https://api2.openreview.net.}}.
Specifically, we crawl NeurIPS 2025 and ICLR 2025 papers from OpenReview, filtering out withdrawn submissions without review comments and placeholder papers lacking full text.
The remaining papers are partitioned into four strata according to their final decisions (Reject, Poster, Spotlight, Oral), from which we perform stratified sampling.
We specially include every submission track in the conference within each stratum to ensure idea diversity.
We apply our extraction agent $\boldsymbol{\mathcal{M}}_e$ to harvest ideas from each paper, followed by manual correction and verification, finally yielding 217 point-wise samples, which we mark as $\boldsymbol{\mathcal{D}}_\textrm{point}$.
The final dataset comprises 136 ICLR 2025 ideas and 81 NeurIPS 2025 ideas, including 138 Reject, 66 Poster, 9 Spotlight, and 4 Oral decisions.
Consequently, Reject, Poster, and Highlight (Spotlight + Oral) ideas account for 61.3\%, 29.3\%, and 9.4\% of all the samples, respectively.
Their detailed track distributions are listed in Tab.\ref{tab:track_distribution}.

\paragraph{Tasks and Metrics.}
We have two classification tasks of different difficulty in this setting:
\textit{i) Binary classification}, predicting whether an idea can be accepted. The labels are Reject and Accept (including Poster, Spotlight, and Oral);
\textit{ii) Three-way classification}, where we group Spotlight and Oral into a single class called Highlight. The labels are Reject, Poster, and Highlight (including Spotlight and Oral).
We use Accuracy and macro F1 score as the evaluation metrics.

\subsection{Group-wise Dataset}
\paragraph{Dataset Construction.}
To construct idea groups with similar topics, we use the abstract of each idea in $\boldsymbol{\mathcal{D}}_\textrm{point}$ as a query to retrieve similar papers from all the papers we crawl from ICLR 2025 and NeurIPS 2025.
We adopt \texttt{bge-base-en-v1.5}\footnote{\url{https://huggingface.co/BAAI/bge-large-en-v1.5}} as the embedding model to retrieve the 800 most similar papers, then apply \texttt{bge-reranker-base}\footnote{\url{https://huggingface.co/BAAI/bge-reranker-base}} as the reranker to reduce the pool to 120.
From these 120 candidates, we select the highest-similarity paper in each label stratum to form a group, enabling automatic ranking of the ideas by their labels.
We apply the same extraction and verification pipeline to the retrieved papers, finally obtaining 172 group-wise instances as $\boldsymbol{\mathcal{D}}_\textrm{group}$.

\paragraph{Tasks and Metrics.}
We define two objectives for group-wise tasks:
\textit{i) Best Selection} to identify the single best idea within each group, and  
\textit{ii) Ranking} to produce the exact full ranking list of all ideas in the group.  
We evaluate the best-idea selection by Accuracy, and assess the ranking task via Longest Increasing Subsequence (LIS) score as well as Accuracy.
The LIS score is computed as follows: let the group size be $n$ and the length of the longest increasing subsequence of the predicted ranking with respect to the gold ranking be $l$.
Then the LIS score is $\frac{l}{n}$.
For example, if the gold ranking is $\{1, 2, 3, 4\}$ and the predicted ranking is $\{2, 1, 3, 4\}$, their longest increasing subsequence is $\{1, 3, 4\}$ or $\{2, 3, 4\}$, all with a length of 3, so the LIS score is $\frac{3}{4}=0.75$.
And for \textit{ranking} tasks, the Accuracy$=1$ if and only if the predicted ranking is identical to the gold ranking, otherwise the Accuracy will be $0$.

\subsection{Pair-wise Dataset}
Since idea pairs naturally exist within each group, our pairwise dataset is constructed on top of $\boldsymbol{\mathcal{D}}_{\textrm{group}}$.  
For a group of size $n$, we can sample $\binom{n}{2}$ pairs. We perform the sampling at two difficulty levels:
1) \textit{easy}: pairs whose two ideas have markedly different labels (Highlight vs. Reject), and are therefore easier to distinguish.  
2) \textit{difficult}: pairs whose two ideas have similar labels (Poster vs. Reject, Highlight vs. Poster), and are therefore harder to distinguish.
After screening and filtering, we construct a dataset $\boldsymbol{\mathcal{D}}_{\textrm{pair}}$ comprising 372 samples, including 172 easy pairs and 200 difficult pairs.
We directly use Accuracy to evaluate pair-wise tasks and report the results of easy and difficult pairs separately.